\section{Background}\label{sec:background}

The gap between a fuzzing campaign and a usable bug report is wider than it should
be, and it is made of several compounding obstacles. Wireless fuzzing leaves the
task of reproduction unfinished. The medium then hands every replay a noisy,
unreliable oracle. And the one classical algorithm that would polish the result
rests on an assumption the noise invalidates. We deal with these in turn, because
RDD is built to resolve them together.

\subsection{Over-the-air fuzzing and the reproduction bottleneck}

Fuzzing wireless stacks keeps revealing real vulnerabilities.
\textsc{SweynTooth}~\cite{sweyntooth2020} found twelve Bluetooth Low Energy flaws
across seven vendors, \textsc{BrakTooth}~\cite{braktooth2022} sixteen in Bluetooth
Classic, and \textsc{U-Fuzz}~\cite{ufuzz2024} carried the same idea to stateful
IoT protocols, 5G\,NR among them, on commercial devices. Each campaign ends with a
large packet trace and a thin crash log, a terse record that a fault occurred but
not which of the many fuzzed packets~\cite{airbugcatcher2024} was responsible. The
work does not scale well. Every distinct crash signature needs its own
reconstruction, and aligning a sparse crash log against a vast trace by hand is a
bottleneck that keeps teams from acting on what their fuzzers already found.

Turning that raw output into something a developer can use means
\emph{reproduction}: reducing the trace to the minimal packet sequence that still
triggers the crash. Naive record-and-replay fails at once, because a wireless
session depends on values that were valid only in the moment, like authentication
challenges, sequence numbers, and freshly negotiated parameters, none of which a
stale replay can satisfy. \textsc{AirBugCatcher}~\cite{airbugcatcher2024} is the
state of the art, and it works in two stages. First, offline, it groups crashes
that share a signature and walks a backward window from each crash to collect the
fuzzed packets most likely responsible, distilling each group into a concise set
of suspect packets. Then, online, it rebuilds the attack on a fresh, live
connection. Acting as a man-in-the-middle, it lets the handshake proceed,
intercepts and mutates only the suspect packets, and re-runs the scenario,
repeating within a time budget until the expected crash returns or the budget is
spent. Two design choices underpin the approach: the suspect set is kept small,
and the medium's non-determinism is handled by repeating trials under a budget
rather than by making any single verdict trustworthy.

\subsection{The oracle problem and flakiness under noise}

In software testing, the \emph{oracle} is the mechanism that decides whether a run
passed or failed. The oracle problem~\cite{oracleproblem2015} is the long-standing
observation that automating this verdict is often harder than building the system
under test. Even on ordinary wired software, verdicts can wobble. \emph{Flaky}
tests, which pass or fail without any code change, are a well-documented nuisance
driven by timing, concurrency, and environment~\cite{flakytests2014}. On a
wireless target the inconsistency becomes physical. Multipath fading, interference,
and radio drift mean identical packets need not yield the same result twice, and
the device's hidden state (timers, queues, the seeds of its pseudo-random
generators) is rarely reset between trials. The component that should answer
``did this recipe reproduce the bug?'' is therefore unreliable by construction.
Repeating a test, as \textsc{AirBugCatcher} does, makes a true reproduction more
likely to show up but never makes the verdict trustworthy. The distinction
matters. A flaky software test can in principle be stabilised by controlling the
timing and the randomness, but radio non-determinism is intrinsic to the medium
and cannot be removed by better test hygiene. The verdict cannot be cleaned. It
has to be \emph{stabilised}.

\subsection{Delta debugging and the monotonicity assumption}

Delta debugging, using the \texttt{ddmin} algorithm~\cite{ddmin2002}, is a precise
tool for input minimisation. It works by splitting the suspect elements, keeping
the half that still causes failure, and repeating this process until removing any
single element stops the failure. This produces a \emph{1-minimal} result, usually
found in $O(\log n)$ tests if only one packet is at fault, or $O(n^2)$ in the
worst case. For example, if a recipe has eight suspect packets and only one is
responsible, \texttt{ddmin} tests the first four. If the crash still happens, it
discards the other four, then narrows down to two, then one, stopping when no more
packets can be removed without losing the crash. Importantly, \texttt{ddmin}
guarantees that every remaining packet is needed for the failure, but it does not
ensure that a smaller, different subset could not also cause the failure. Finding
the absolute minimum is NP-complete~\cite{hdd2006}, so practical tools rely on
this local guarantee.

These methods depend on certain assumptions about the test results, as explained
in the original work. \texttt{ddmin} requires the oracle to be \emph{consistent},
meaning the same input always gives the same result, and it assumes failures are
\emph{monotone}, so if a set of packets causes a crash, adding more packets will
not stop it. Channel noise breaks the consistency assumption. If a test fails to
reproduce the crash, the algorithm cannot tell if it was due to an innocent half
or just bad luck with the radio, so it might discard the half with the real
culprit and end up with the wrong answer or no answer at all. The monotonicity
assumption is also weak: a single \emph{suppressor} packet can hide a crash that a
smaller subset would still cause. We discuss this harder case in
Section~\ref{sec:our_approach}.

\subsection{Where that leaves \textsc{AirBugCatcher}}

Because it cannot rely on the oracle, \textsc{AirBugCatcher} switches from delta
debugging to bounded enumeration. It tests combinations of suspect packets up to a
set size and repeats each test within a time limit. This approach is practical but
has two main drawbacks. First, the size limit means that if the real recipe is
larger, the tool will never find it, no matter how long it runs. Second, repeating
tests increases the chance of seeing a real reproduction but does not make any
single result reliable. \textsc{AirBugCatcher} checks each crash against the
expected signature before accepting it, flagging mismatches instead of counting
them. However, this exact match is not perfect on closed firmware, so a crash with
the same signature as a different fault can still be accepted. As a result, two
types of errors build up: real reproductions lost to channel noise, and
\emph{false credits}, where PoCs are reported as working but do not actually
reproduce the bug. To bring back delta debugging, we need to restore a trustworthy
verdict despite the noise, which is the focus of the next section.
